\documentclass{article}
\usepackage{mathtools} 
\usepackage{fontspec}
\usepackage[UTF8]{ctex}
\usepackage{amsthm}
\usepackage{mdframed}
\usepackage{xcolor}
\usepackage{amssymb}
\usepackage{amsmath}


% 定义新的带灰色背景的说明环境 zremark
\newmdtheoremenv[
  backgroundcolor=gray!10,
  % 边框与背景一致，边框线会消失
  linecolor=gray!10
]{zremark}{说明}

% 通用矩阵命令: \flexmatrix{矩阵名}{元素符号}{行数}{列数}
\newcommand{\flexmatrix}[4]{
  \[
  #1 = \begin{pmatrix}
    #2_{11}     & #2_{12}     & \cdots & #2_{1#4}   \\
    #2_{21}     & #2_{22}     & \cdots & #2_{2#4}   \\
    \vdots      & \vdots      & \ddots & \vdots     \\
    #2_{#31}    & #2_{#32}    & \cdots & #2_{#3#4}
  \end{pmatrix}
  \]
}

% 简化版命令（默认矩阵名为A，元素符号为a）: \quickmatrix{行数}{列数}
\newcommand{\quickmatrix}[2]{\flexmatrix{A}{a}{#1}{#2}}

\begin{document}
\title{注释 6.4}
\author{张志聪}
\maketitle

\begin{zremark}
  任意两点$x_1 < x_2$，对任意$\lambda \in (0, 1)$，
  $y = \lambda x_1 + (1 - \lambda)x_2$，
  则有，$y \in (x_1, x_2)$。
\end{zremark}
\textbf{证明：}
\begin{align*}
  y - x_1
   & = \lambda x_1 + (1 - \lambda)x_2 - x_1  \\
   & = (\lambda - 1) x_1 + (1 - \lambda) x_2 \\
   & = (1 - \lambda) (x_2 - x_1)             \\
   & > 0
\end{align*}
又
\begin{align*}
  y - x_2
   & = \lambda x_1 + (1 - \lambda)x_2 - x_2 \\
   & = \lambda x_1 - \lambda x_2            \\
   & = \lambda (x_1 - x_2)                  \\
   & < 0
\end{align*}
$\therefore y \in (x_1, x_2)$.

\begin{zremark}
  定义1 中不等式两边的几何含义。
\end{zremark}
\textbf{证明：}

以凸函数为例，任取$(x_1, x_2)$（设$x_1 < x_2$）中的一点$x' = \lambda x_1 + (1 - \lambda)x_2$，
对应的函数值$f(x') = f(\lambda x_1 + (1 - \lambda)x_2)$。

点$(x_1, f(x_1)), (x_2, f(x_2))$确定直线：
\begin{align*}
  \frac{y - y_1}{y_2 - y_1}          & = \frac{x - x_1}{x_2 - x_1}                            \\
  \frac{y - f(x_1)}{f(x_2) - f(x_1)} & = \frac{x - x_1}{x_2 - x_1}                            \\
  y                                  & = [f(x_2) - f(x_1)] \frac{x - x_1}{x_2 - x_1} + f(x_1)
\end{align*}
于是，$x = x'$时，我们有
\begin{align*}
   & [f(x_2) - f(x_1)] \frac{\lambda x_1 + (1 - \lambda)x_2 - x_1}{x_2 - x_1} + f(x_1) \\
   & = [f(x_2) - f(x_1)] \frac{(1 - \lambda)(x_2 - x_1)}{x_2 - x_1} + f(x_1)           \\
   & = [f(x_2) - f(x_1)](1 - \lambda) + f(x_1)                                         \\
   & = (\lambda - 1)f(x_1) + (1 - \lambda)f(x_2) + f(x_1)                              \\
   & = \lambda f(x_1) + (1 - \lambda)f(x_2)
\end{align*}

综上，函数在$f(x_1), f(x_2)$两点处的形成的割线在$x \in (x_1, x_2)$上与函数值$f(x)$的比较，
定义出凸函数和凹函数。

\begin{zremark}
  $f$为$I$上的凸函数的充分条件是：对于$I$上任意三个点$x_1 < x_2 < x_3$，有
  \begin{align*}
    \frac{f(x_2) - f(x_1)}{x_2 - x_1} \leq \frac{f(x_3) - f(x_1)}{x_3 - x_1} \leq \frac{f(x_3) - f(x_2)}{x_3 - x_2}
  \end{align*}
\end{zremark}

\textbf{证明：}

引理保证了前一个不等式，我们只需证明后一个不等式。

\begin{itemize}
  \item 必要性

        在书中已有
        \begin{align*}
          (x_3 - x_1) f(x_2) \leq (x_3 - x_2)f(x_1) + (x_2 - x_1)f(x_3)
        \end{align*}
        我们在此基础上整理：
        \begin{align*}
          (x_3 - x_1) f(x_2) \leq (x_3 - x_2)f(x_1) + (x_2 + x_3 - x_3 - x_1)f(x_3)          \\
          (x_3 - x_1) f(x_2) \leq (x_3 - x_2)f(x_1) + (x_3 - x_1)f(x_3) + (x_2 - x_3)f(x_3)  \\
          (x_3 - x_2)f(x_3) - (x_3 - x_2)f(x_1) \leq  (x_3 - x_1)f(x_3) - (x_3 - x_1) f(x_2) \\
          \frac{f(x_3) - f(x_1)}{x_3 - x_1} \leq \frac{f(x_3 - f(x_2))}{x_3 - x_2}
        \end{align*}

  \item 充分性

        和书中类似

\end{itemize}

\end{document}